\chapter{Metodologia}
\label{cap:metodologia}

Este capítulo descreve a contribuição do trabalho quanto à extensão do MellowDB para ambientes distribuídos e o método experimental usado para caracterizá-la. As seções a seguir tratam da engenharia em volta da visão geral da extensão, a adaptação a \textit{replica sets} e \textit{clusters} particionados e o arcabouço que torna os experimentos reprodutíveis. A última seção define o método experimental dos conjuntos de dados, design das séries de \textit{benchmark} e métricas. Os resultados numéricos ficam reservados para o Capítulo~\ref{cap:resultados}.

\section{Visão geral da extensão}
\label{sec:visao-geral}

O capítulo anterior mostrou que a evolução semântica dos dados é um problema de significado, e não de esquema, e que o MellowDB (Seção~\ref{sec:mellowdb}) o resolve por meio de uma arquitetura de nó único de MongoDB, através das estratégias \textit{rewrite} e \textit{preprocess}. Entretanto, isso carrega uma limitação evidente ao analisar bancos que perduram por longos períodos, como o já mencionado DATASUS, que para a pesquisa de Nepomuceno contava com cerca de 6,1 milhões de registros e 220 operações de evolução. Um volume de dados muito alto pode prejudicar a performance de consultas quando restrito a um só servidor.

A solução é aplicar a distribuição dos dados focando nos pontos levantados na Seção~\ref{sec:distribuidos}. De um lado, replicar os dados para garantir a tolerância a falhas e distribuir a carga de leitura e por outro, repartir a coleção para evitar que um servidor seja sobrecarregado de informação. Este é um campo não muito explorado para a evolução semântica e que é sugerido como trabalho futuro na tese de Nepomuceno, tornando-se o objetivo central desta contribuição e das decisões que se seguem.

O MellowDB foi desenvolvido sobre o MongoDB e, além das conveniências de um banco de documentos \textit{schemaless} \citep{cattell2011} para as operações de evolução, pode-se contar com os mecanismos já existentes de distribuição para Python, linguagem de programação em que o sistema original\footnote{Repositório da implementação de Nepomuceno: \url{https://github.com/pisn/semantic_heterogeneous_database}} e esta extensão\footnote{Repositório da contribuição deste trabalho: \url{https://github.com/gust-vaz/semantic_heterogeneous_extended}} foram contruídos. A decisão de manter o mesmo SGBD advém dessa herança e praticidade, mas também implica em limitações das primitivas que o próprio MongoDB disponibiliza.

A outra decisão foi sobre como expor a extensão a quem a usa. Buscou-se preservar o modelo de programação original em que a mesma classe \texttt{BasicCollection} que servia um nó isolado deveria servir, sem reescrita de código, as abordagens distribuídas. O que varia de uma implementação para outra reduz-se à \textit{connection string} e aos três parâmetros \texttt{write\_concern}, o par \texttt{read\_uri}/\texttt{read\_mode} e \texttt{shard\_key}. Para que isso funcionasse sem configuração manual, a coleção passou a detectar a topologia em tempo de execução, ou seja, ao ser inicializada, checa o sistema, já operando, com o comando \texttt{hello} e decide como agir:
\begin{itemize}
    \item Um nome de conjunto (\texttt{setName}) revela um \textit{replica set} e ativa as transações da cadeia de versões;
    \item A resposta \texttt{msg} igual a \texttt{isdbgrid} revela um roteador \texttt{mongos} e ativa o particionamento;
    \item Um \textit{mongod} isolado não dispara nenhum dos dois e preserva o comportamento herdado.
\end{itemize}

Essas duas decisões precedem o que as próximas seções detalham, sempre dentro das restrições herdadas do MongoDB. Preservando os mecanismos descritos na Seção~\ref{sec:mongodb}, privilegia-se as implicações das escolhas de implementação desse projeto de extensão.

\section{Extensão para \textit{replica sets}}
\label{sec:ext-replicaset}

Levar o MellowDB à forma mais simples de distribuição, a replicação, se encontra com a cadeia de versões, um ponto sensível do funcionamento do sistema. Ela é a fonte de verdade que traduz cada consulta entre termos de épocas diferentes, de modo que, se for escrita pela metade ou lida de uma cópia desatualizada, a tradução passa a devolver resultados errados. Por isso foi importante considerar como armazená-la com segurança e de onde lê-la.

Aplicar uma operação semântica compreende alguns passos, pois atualiza-se o nó terminal da cadeia, com os ponteiros \texttt{next\_operation} e \texttt{next\_version}, e insere-se o novo nó que passa a descrever a versão atual. Em um único servidor, essa sequência ou falhava por inteiro ou era refeita, porém, em um \textit{replica set}, uma eleição ou a queda do primário no meio dela poderia deixar a cadeia em um estado intermediário, com um ponteiro já gravado e o nó correspondente ausente. Assim, quando a coleção detecta que opera sobre um \textit{replica set}, a escrita da cadeia é envolvida em uma transação multidocumento, confirmada com \textit{write concern} \texttt{majority} e lida sob \textit{read concern} \texttt{snapshot}, uma forma de indicar que ou, a operação inteira é aceita por uma maioria dos nós, e sobrevive a qualquer eleição posterior, ou nenhuma parte dela é aplicada. No caso de operações semânticas simultâneas, aplica-se um controle de concorrência otimista, onde, ao invés de bloquear a cadeia previamente, detecta-se e rejeita-se a escrita que chegou por último, lançando uma exceção de conflito (\texttt{VersionChainConflict}).

A durabilidade das demais escritas, como os registros brutos e os processados, é exposta como o parâmetro \texttt{write\_concern}, que o usuário ajusta conforme o compromisso desejado entre segurança e latência, de \texttt{w:1}, que confirma assim que o primário grava, a \texttt{w:majority}, que aguarda a maioria dos nós. O custo dessa garantia, à medida que o conjunto de nós cresce, é um dos pontos que os testes mais adiante procuram caracterizar.

Do lado da leitura, o retorno da replicação é poder consultar a partir de qualquer réplica em vez de apenas no primário. A extensão aproveita isso preservando a distinção entre os dois tipos de leitura que realiza, a da cadeia de versões, que idealmente precisa estar sempre atualizada, e a dos registros, que pode tolerar algum grau de defasagem. Esse comportamento é controlado por \texttt{read\_uri}, que indica o nó de onde os registros são lidos, e por \texttt{read\_mode}, com dois modos:
\begin{itemize}
    \item \texttt{split} (padrão): a cadeia de versões é sempre lida do primário, sob \textit{read concern} \texttt{majority}, enquanto apenas os registros, a maior carga de dados, são recuperados pelo nó indicado em \texttt{read\_uri}. O volume é distribuido sem arriscar a correção da tradução;
    \item \texttt{single\_source}: tudo, cadeia e registros, é lido de um mesmo nó, assumindo de propósito a defasagem. Assim, servindo para ler de forma consistente de um nó específico ou para medir o efeito dessa defasagem.
\end{itemize}

O nó e o modo de leitura podem ser trocados em tempo de execução, sem reconstruir a coleção, o que permite redirecionar as consultas para outro secundário no meio de um experimento.

Em suma, a extensão para \textit{replica sets} mantém a cadeia de versões autoritativa no primário e oferece a replicação como ganho de durabilidade e de vazão de leitura, sem comprometer a correção da tradução semântica.

%\section{Extensão para \textit{sharding}}
%\label{sec:ext-sharding}

%\section{Ferramental, infraestrutura e reprodutibilidade}
%\label{sec:ferramental}

%\section{Método experimental}
%\label{sec:metodo}